\documentclass[11pt,a4paper]{amsart}
\usepackage[T1]{fontenc}
\usepackage{lmodern}
\usepackage[brazil]{babel}
\usepackage{amsmath,amssymb,amsfonts,amsthm}
\usepackage{enumitem}
\usepackage{microtype}
\usepackage[left=2.3cm,right=2.3cm,top=2.6cm,bottom=2.6cm]{geometry}

%==================== Comandos ====================
\newcommand{\R}{\mathbb{R}}
\newcommand{\Z}{\mathbb{Z}}
\newcommand{\N}{\mathbb{N}}
\renewcommand{\d}{\mathrm{d}}
\newcommand{\Q}{\mathcal{Q}}
\newcommand{\Adm}{\mathcal{A}}

%==================== Ambientes ====================
\theoremstyle{plain}
\newtheorem{teorema}{Teorema}[section]
\newtheorem{lema}[teorema]{Lema}
\newtheorem{proposicao}[teorema]{Proposição}
\newtheorem{corolario}[teorema]{Corolário}

\theoremstyle{definition}
\newtheorem{definicao}[teorema]{Definição}
\newtheorem{exercicio}[teorema]{Exercício}

\theoremstyle{remark}
\newtheorem{observacao}[teorema]{Observação}
\newtheorem{pergunta}[teorema]{Pergunta}

%==================== Título ====================
\title[A desigualdade de Wirtinger]{A desigualdade de Wirtinger e a desigualdade
isoperim\'etrica no plano}
\author{}
\date{}

\begin{document}

\begin{abstract}
Estas notas provam a desigualdade de Wirtinger unidimensional (tamb\'em chamada
de desigualdade de Poincar\'e--Wirtinger peri\'odica) por um argumento
variacional t\'ipico de um curso de equa\c{c}\~oes diferenciais ordin\'arias:
sup\~oe-se a exist\^encia de um extremizador, deriva-se a equa\c{c}\~ao de
Euler--Lagrange associada, resolve-se a equa\c{c}\~ao diferencial de segunda
ordem resultante e, a partir das solu\c{c}\~oes, l\^e-se a constante \'otima.
Em nenhum momento usamos s\'eries de Fourier. Em seguida aplicamos a
desigualdade \`a prova de Hurwitz da desigualdade isoperim\'etrica no plano.
\end{abstract}

\maketitle

{\small\noindent Notas de aula para um curso de Equa\c{c}\~oes Diferenciais
Ordin\'arias (segundo ano). O p\'ublico-alvo s\~ao estudantes que j\'a
conhecem c\'alculo de v\'arias vari\'aveis e a teoria elementar de EDOs
lineares de segunda ordem.}

\vspace{3mm}

%=====================================================================
\section{Introdu\c{c}\~ao}
%=====================================================================

Comparar o ``tamanho'' de uma fun\c{c}\~ao com o ``tamanho'' de sua derivada
\'e uma das ideias mais f\'erteis da an\'alise. A desigualdade de Wirtinger \'e
o prot\'otipo unidimensional dessa compara\c{c}\~ao: ela afirma que, para uma
fun\c{c}\~ao peri\'odica de m\'edia nula, a energia $\int u^2$ nunca excede a
energia da derivada $\int (u')^2$. Intuitivamente, oscilar (ter m\'edia nula)
custa derivada, e quanto mais devagar a fun\c{c}\~ao oscila, mais pr\'oxima a
energia da fun\c{c}\~ao fica da energia da derivada.

O ponto de vista que adotaremos \'e o do \emph{c\'alculo de varia\c{c}\~oes}.
Em vez de expandir $u$ em uma s\'erie de senos e cossenos, vamos procurar a
fun\c{c}\~ao que torna a raz\~ao $\int u^2 / \int (u')^2$ a maior poss\'\i vel.
Tal fun\c{c}\~ao extremal, se existir, deve satisfazer uma equa\c{c}\~ao
diferencial --- a \emph{equa\c{c}\~ao de Euler--Lagrange} --- que neste caso \'e
simplesmente a equa\c{c}\~ao do oscilador harm\^onico
\[
  u'' + \omega^2 u = 0 .
\]
Resolver essa EDO e impor periodicidade revela que os \'unicos extremais
poss\'\i veis s\~ao os ``harm\^onicos'' $A\cos kt + B\sin kt$, e a melhor
escolha \'e o primeiro harm\^onico $k=1$. \'E dessa an\'alise de EDO que a
constante \'otima emerge.

Na \'ultima se\c{c}\~ao colhemos o fruto geom\'etrico: a desigualdade de
Wirtinger fornece, quase que de imediato, a desigualdade isoperim\'etrica no
plano,
\[
  4\pi A \le L^2 ,
\]
que diz que, dentre todas as curvas fechadas de comprimento $L$ fixado, o
c\'\i rculo \'e a que encerra a maior \'area. Essa dedu\c{c}\~ao \'e devida a
Hurwitz.

%=====================================================================
\section{A classe admiss\'\i vel e o enunciado}
%=====================================================================

Fixamos o per\'\i odo $2\pi$ ao longo de quase todo o texto; a forma geral, com
per\'\i odo $L$, ser\'a obtida por uma mudan\c{c}a de escala na
Se\c{c}\~ao~\ref{sec:escala}.

\begin{definicao}[Classe admiss\'\i vel]\label{def:adm}
Seja $\Adm$ o conjunto das fun\c{c}\~oes $u\colon\R\to\R$ de classe $C^1$,
peri\'odicas de per\'\i odo $2\pi$ (isto \'e, $u(t+2\pi)=u(t)$ para todo $t$) e
de \emph{m\'edia nula}:
\[
  \int_0^{2\pi} u(t)\,\d t = 0 .
\]
\end{definicao}

Observe que se $u$ \'e $C^1$ e $2\pi$-peri\'odica, ent\~ao $u'$ tamb\'em \'e
$2\pi$-peri\'odica e cont\'\i nua, de modo que todas as integrais que
escreveremos s\~ao finitas.

\begin{teorema}[Desigualdade de Wirtinger]\label{thm:wirtinger}
Para toda $u\in\Adm$ vale
\[
  \int_0^{2\pi} u(t)^2 \,\d t \;\le\; \int_0^{2\pi} u'(t)^2 \,\d t .
\]
Al\'em disso, a igualdade ocorre \emph{se, e somente se,}
\[
  u(t) = a\cos t + b\sin t
\]
para constantes $a,b\in\R$.
\end{teorema}

A constante $1$ que multiplica $\int (u')^2$ \'e \emph{\'otima}: nenhuma
constante menor que $1$ faria a desigualdade valer para toda $u\in\Adm$, pois
ela j\'a se torna uma igualdade nos primeiros harm\^onicos. Demonstrar este
teorema --- incluindo a otimalidade da constante e a descri\c{c}\~ao completa
do caso de igualdade --- \'e o objetivo das Se\c{c}\~oes
\ref{sec:rayleigh}--\ref{sec:conclusao}.

\begin{pergunta}\label{perg:media}
\textbf{Por que a condi\c{c}\~ao de m\'edia nula \'e indispens\'avel?}
Sem ela a desigualdade \'e falsa. De fato, tome a fun\c{c}\~ao constante
$u\equiv c\neq 0$. Ent\~ao $u'\equiv 0$, logo
\[
  \int_0^{2\pi} u^2 = 2\pi c^2 > 0 = \int_0^{2\pi}(u')^2 ,
\]
e a desigualdade $\int u^2\le\int (u')^2$ falha gritantemente. As constantes
n\~ao nulas s\~ao exatamente as fun\c{c}\~oes peri\'odicas com derivada nula;
a hip\'otese de m\'edia nula serve precisamente para exclu\'\i-las da
competi\c{c}\~ao. Em termos das oscila\c{c}\~oes: $\int u^2$ mede o tamanho de
$u$, mas a derivada s\'o ``enxerga'' a parte oscilante de $u$; impor m\'edia
nula \'e descartar a parte constante, que a derivada ignora.
\end{pergunta}

%=====================================================================
\section{O quociente de Rayleigh e a estrat\'egia variacional}
\label{sec:rayleigh}
%=====================================================================

Para uma fun\c{c}\~ao $u\in\Adm$ \emph{n\~ao identicamente nula}, definimos o
\emph{quociente de Rayleigh}
\begin{equation}\label{eq:rayleigh}
  \Q[u] \;=\; \frac{\displaystyle\int_0^{2\pi} u(t)^2\,\d t}
                   {\displaystyle\int_0^{2\pi} u'(t)^2\,\d t}.
\end{equation}
O denominador \'e estritamente positivo: se fosse $\int (u')^2=0$, ent\~ao
$u'\equiv 0$ (pois $u'$ \'e cont\'\i nua), logo $u$ seria constante; mas a
\'unica constante de m\'edia nula \'e $u\equiv 0$, exclu\'\i da por
hip\'otese. Assim $\Q[u]$ est\'a bem definido em $\Adm\setminus\{0\}$.

A desigualdade de Wirtinger equivale exatamente \`a afirma\c{c}\~ao
\[
  \sup_{u\in\Adm\setminus\{0\}} \Q[u] \;=\; 1 ,
\]
com o supremo sendo atingido. De fato, $\Q[u]\le 1$ para toda $u$ \'e o mesmo
que $\int u^2\le \int (u')^2$, e a igualdade $\Q[u]=1$ corresponde ao caso de
igualdade do Teorema~\ref{thm:wirtinger}.

\medskip

\noindent\textbf{A hip\'otese de exist\^encia.}\quad
A estrat\'egia variacional consiste em estudar um \emph{maximizador} de $\Q$.
Para um curso de EDOs, faremos a seguinte hip\'otese de trabalho, e o leitor
deve t\^e-la sempre em mente:

\begin{observacao}[Caveat anal\'\i tico]\label{obs:existencia}
\emph{Admitiremos que o supremo de $\Q$ em $\Adm\setminus\{0\}$ \'e atingido},
isto \'e, que existe $u_\star\in\Adm\setminus\{0\}$ com
$\Q[u_\star]=\sup_u\Q[u]$. Justificar rigorosamente essa exist\^encia exige o
m\'etodo direto do c\'alculo de varia\c{c}\~oes em espa\c{c}os de Sobolev
(compacidade da inclus\~ao $H^1\hookrightarrow L^2$), o que est\'a fora do
escopo destas notas. Admitiremos ainda que o maximizador \'e suficientemente
regular ($C^2$); a teoria de regularidade garante isso, mas, como veremos na
Observa\c{c}\~ao~\ref{obs:bootstrap}, a pr\'opria equa\c{c}\~ao de
Euler--Lagrange j\'a indica por que a regularidade \'e autom\'atica.
\end{observacao}

A partir de agora, $u$ denota um maximizador fixado. Como $\Q$ \'e
invariante por reescala ($\Q[cu]=\Q[u]$ para $c\neq 0$), podemos
\emph{normalizar} $u$ de modo que
\begin{equation}\label{eq:normaliza}
  \int_0^{2\pi} u'(t)^2\,\d t = 1 .
\end{equation}
Com essa normaliza\c{c}\~ao, maximizar $\Q[u]$ \'e o mesmo que maximizar
$\int_0^{2\pi}u^2$ sujeito \`as duas restri\c{c}\~oes
\begin{equation}\label{eq:restricoes}
  \int_0^{2\pi}(u')^2 = 1 ,
  \qquad
  \int_0^{2\pi} u = 0 .
\end{equation}

%=====================================================================
\section{A equa\c{c}\~ao de Euler--Lagrange}
\label{sec:EL}
%=====================================================================

Vamos extrair de $u$ uma equa\c{c}\~ao diferencial. A ferramenta \'e o
\emph{m\'etodo dos multiplicadores de Lagrange}, aplicado ao problema com
v\'\i nculos \eqref{eq:restricoes}.

\subsection{A condi\c{c}\~ao de estacionariedade}

Como $u$ maximiza $\int u^2$ sob os v\'\i nculos \eqref{eq:restricoes}, o
m\'etodo dos multiplicadores de Lagrange diz que existem constantes
$\lambda,\mu\in\R$ tais que $u$ \'e um ponto cr\'\i tico \emph{livre} do
funcional
\begin{equation}\label{eq:lagrangiano}
  \mathcal{L}[w]
  \;=\; \int_0^{2\pi} w^2
        \;-\;\lambda\!\left(\int_0^{2\pi}(w')^2 - 1\right)
        \;-\;\mu\int_0^{2\pi} w .
\end{equation}
``Ponto cr\'\i tico livre'' significa que, para \emph{qualquer} fun\c{c}\~ao
de teste $v$ de classe $C^1$ e $2\pi$-peri\'odica (sem nenhuma
restri\c{c}\~ao sobre $v$), a fun\c{c}\~ao de uma vari\'avel
$\varepsilon\mapsto \mathcal{L}[u+\varepsilon v]$ tem derivada nula em
$\varepsilon=0$.

Calculemos essa derivada. Expandindo,
\[
  \mathcal{L}[u+\varepsilon v]
  = \int (u+\varepsilon v)^2
    - \lambda\!\left(\int (u'+\varepsilon v')^2 - 1\right)
    - \mu\int (u+\varepsilon v),
\]
e derivando em $\varepsilon$ no ponto $\varepsilon=0$ obtemos a
\emph{primeira varia\c{c}\~ao}
\begin{equation}\label{eq:primvar}
  \left.\frac{\d}{\d\varepsilon}\right|_{\varepsilon=0}\!\!\mathcal{L}[u+\varepsilon v]
  = 2\int_0^{2\pi} u\,v
    - 2\lambda\int_0^{2\pi} u'\,v'
    - \mu\int_0^{2\pi} v .
\end{equation}
A condi\c{c}\~ao de estacionariedade \'e que \eqref{eq:primvar} se anule para
toda fun\c{c}\~ao de teste peri\'odica $v$:
\begin{equation}\label{eq:fraca}
  2\int_0^{2\pi} u\,v
  - 2\lambda\int_0^{2\pi} u'\,v'
  - \mu\int_0^{2\pi} v = 0
  \qquad\text{para toda $v\in C^1$, $2\pi$-peri\'odica.}
\end{equation}

\subsection{Integra\c{c}\~ao por partes e o termo de fronteira}

O \'unico termo de \eqref{eq:fraca} que ainda cont\'em uma derivada de $v$ \'e
$\int u'v'$. Integramos por partes para passar a derivada de $v$ para $u$.
Lembrando que admitimos $u\in C^2$ (Observa\c{c}\~ao~\ref{obs:existencia}),
\begin{equation}\label{eq:partes}
  \int_0^{2\pi} u'\,v' \,\d t
  = \underbrace{\big[\,u'(t)\,v(t)\,\big]_0^{2\pi}}_{=\,0}
    - \int_0^{2\pi} u''\,v \,\d t
  = - \int_0^{2\pi} u''\,v \,\d t .
\end{equation}

\begin{pergunta}\label{perg:fronteira}
\textbf{Onde entram as condi\c{c}\~oes de contorno peri\'odicas?}
Exatamente aqui. O termo de fronteira \'e
$u'(2\pi)v(2\pi) - u'(0)v(0)$. Como $u'$ e $v$ s\~ao $2\pi$-peri\'odicas,
temos $u'(2\pi)=u'(0)$ e $v(2\pi)=v(0)$, e o termo se cancela
\emph{automaticamente}. \'E isto que torna o problema peri\'odico t\~ao limpo:
n\~ao h\'a fronteira ``de verdade'', e portanto nenhuma condi\c{c}\~ao de
contorno extra precisa ser imposta sobre $v$ --- ela \'e arbitr\'aria entre as
fun\c{c}\~oes peri\'odicas. (Em um problema de Dirichlet com extremos fixos,
ao contr\'ario, exigir\'\i amos $v(0)=v(2\pi)=0$ para anular o termo de
fronteira.)
\end{pergunta}

Substituindo \eqref{eq:partes} em \eqref{eq:fraca} e agrupando,
\begin{equation}\label{eq:integral-nula}
  \int_0^{2\pi} \big(\,2u + 2\lambda u'' - \mu\,\big)\,v\,\d t = 0
  \qquad\text{para toda $v\in C^1$, $2\pi$-peri\'odica.}
\end{equation}

\subsection{Do integral ao pontual: o lema fundamental}

Para converter \eqref{eq:integral-nula} numa equa\c{c}\~ao diferencial,
usamos o lema fundamental do c\'alculo de varia\c{c}\~oes.

\begin{lema}[Lema fundamental do c\'alculo de varia\c{c}\~oes]\label{lema:fundamental}
Seja $h\colon\R\to\R$ cont\'\i nua e $2\pi$-peri\'odica. Se
\[
  \int_0^{2\pi} h(t)\,v(t)\,\d t = 0
\]
para toda fun\c{c}\~ao $v$ de classe $C^1$ e $2\pi$-peri\'odica, ent\~ao
$h\equiv 0$.
\end{lema}

\begin{proof}
Suponha, por absurdo, que $h(t_0)\neq 0$ para algum $t_0$; sem perda de
generalidade $h(t_0)>0$ (caso contr\'ario troque $h$ por $-h$). Por
continuidade, existe $\delta>0$ tal que $h(t) > \tfrac12 h(t_0) > 0$ para todo
$t$ no intervalo $I=(t_0-\delta,\,t_0+\delta)$, que podemos supor contido num
\'unico per\'\i odo. Escolha uma fun\c{c}\~ao ``bolha'' (bump)
$v\colon\R\to\R$ de classe $C^1$, $2\pi$-peri\'odica, com $v\ge 0$,
$v\equiv 0$ fora de $I$ (a menos de transla\c{c}\~oes por $2\pi$) e
$v(t_0)>0$. Por exemplo, pode-se tomar a extens\~ao peri\'odica de
$v(t)=\big(\delta^2-(t-t_0)^2\big)^2$ em $I$ e $0$ fora, que \'e $C^1$. Ent\~ao
\[
  \int_0^{2\pi} h\,v \,\d t = \int_I h\,v\,\d t
  \;\ge\; \tfrac12 h(t_0)\int_I v\,\d t \;>\; 0 ,
\]
contradizendo a hip\'otese. Logo $h\equiv 0$.
\end{proof}

Aplicando o Lema~\ref{lema:fundamental} \`a fun\c{c}\~ao cont\'\i nua e
peri\'odica $h = 2u + 2\lambda u'' - \mu$ em \eqref{eq:integral-nula}, obtemos
a equa\c{c}\~ao \emph{pontual}
\begin{equation}\label{eq:EL}
  2u + 2\lambda u'' - \mu = 0,
  \qquad\text{ou seja}\qquad
  u + \lambda u'' = \frac{\mu}{2}.
\end{equation}
Esta \'e a equa\c{c}\~ao de Euler--Lagrange do nosso problema.

\begin{observacao}[Por que a regularidade \'e ``de gra\c{c}a'']\label{obs:bootstrap}
A identidade \eqref{eq:EL} pode ser obtida supondo apenas $u\in C^1$ (na forma
fraca $2\lambda u'' = \mu - 2u$ no sentido das distribui\c{c}\~oes). Mas, uma
vez sabido que $\lambda\neq 0$ (o que verificaremos abaixo), ela diz que
$u'' = (\mu - 2u)/(2\lambda)$ \'e uma fun\c{c}\~ao \emph{cont\'\i nua} (pois
$u$ \'e cont\'\i nua); logo $u\in C^2$. E ent\~ao o lado direito \'e $C^1$, de
modo que $u\in C^3$, e assim por diante: $u\in C^\infty$. Este \'e o fen\^omeno
de \emph{bootstrap} de regularidade, e \'e por isso que pudemos, sem grande
abuso, admitir $u\in C^2$ desde o in\'\i cio.
\end{observacao}

\subsection{Eliminando o multiplicador $\mu$}

Integramos \eqref{eq:EL} sobre um per\'\i odo:
\[
  \int_0^{2\pi} u \,\d t
  + \lambda \int_0^{2\pi} u'' \,\d t
  = \frac{\mu}{2}\int_0^{2\pi}\d t = \mu\pi .
\]
Os dois termos do lado esquerdo se anulam:
\[
  \int_0^{2\pi} u \,\d t = 0
  \quad\text{(m\'edia nula, Defini\c{c}\~ao~\ref{def:adm});}
  \qquad
  \int_0^{2\pi} u'' \,\d t = \big[\,u'(t)\,\big]_0^{2\pi} = 0
  \quad\text{(periodicidade de $u'$).}
\]
Logo $\mu\pi = 0$, donde
\begin{equation}\label{eq:mu-zero}
  \mu = 0 .
\end{equation}
A equa\c{c}\~ao de Euler--Lagrange reduz-se assim \`a equa\c{c}\~ao homog\^enea
\begin{equation}\label{eq:EL-homog}
  u + \lambda u'' = 0 .
\end{equation}

\subsection{O sinal do multiplicador $\lambda$}

Resta determinar $\lambda$. Multiplicamos \eqref{eq:EL-homog} por $u$ e
integramos:
\[
  \int_0^{2\pi} u^2 + \lambda\int_0^{2\pi} u\,u'' = 0 .
\]
Integrando por partes o segundo termo (com termo de fronteira nulo, como em
\eqref{eq:partes}),
\[
  \int_0^{2\pi} u\,u'' = -\int_0^{2\pi}(u')^2 = -1
\]
pela normaliza\c{c}\~ao \eqref{eq:normaliza}. Portanto
\begin{equation}\label{eq:lambda}
  \int_0^{2\pi} u^2 - \lambda = 0
  \qquad\Longrightarrow\qquad
  \lambda = \int_0^{2\pi} u^2 = \Q[u] > 0 ,
\end{equation}
onde usamos que $u\not\equiv 0$ d\'a $\int u^2>0$. Em particular,
\emph{$\lambda$ \'e exatamente o valor do quociente de Rayleigh no
maximizador}, e \'e estritamente positivo. Isto nos permite escrever
$\lambda = 1/\omega^2$ com $\omega>0$, e \eqref{eq:EL-homog} torna-se a
equa\c{c}\~ao do oscilador harm\^onico
\begin{equation}\label{eq:oscilador}
  u'' + \omega^2 u = 0,
  \qquad \omega = \frac{1}{\sqrt{\lambda}} > 0 .
\end{equation}

%=====================================================================
\section{Resolu\c{c}\~ao da EDO e os extremais}
\label{sec:edo}
%=====================================================================

A equa\c{c}\~ao \eqref{eq:oscilador} \'e uma EDO linear de segunda ordem com
coeficientes constantes, cujo polin\^omio caracter\'\i stico $r^2+\omega^2=0$
tem ra\'\i zes $r=\pm i\omega$. Logo sua solu\c{c}\~ao geral \'e
\begin{equation}\label{eq:solgeral}
  u(t) = A\cos(\omega t) + B\sin(\omega t),
  \qquad A,B\in\R .
\end{equation}

\subsection{A periodicidade for\c{c}a a frequ\^encia a ser inteira}

A fun\c{c}\~ao \eqref{eq:solgeral} \'e solu\c{c}\~ao da EDO para qualquer
$\omega>0$, mas ela ainda precisa pertencer \`a classe admiss\'\i vel
$\Adm$ --- em particular, ser $2\pi$-peri\'odica.

\begin{lema}\label{lema:freq}
Seja $u(t)=A\cos(\omega t)+B\sin(\omega t)$ com $(A,B)\neq(0,0)$ e $\omega>0$.
Ent\~ao $u$ \'e $2\pi$-peri\'odica se, e somente se, $\omega$ \'e um inteiro
positivo.
\end{lema}

\begin{proof}
Se $\omega=k\in\N$, ent\~ao $\cos(kt)$ e $\sin(kt)$ t\^em per\'\i odo
$2\pi/k$, que divide $2\pi$; logo ambas s\~ao $2\pi$-peri\'odicas, e $u$
tamb\'em. Reciprocamente, suponha $u$ $2\pi$-peri\'odica. Podemos escrever
$u(t)=R\cos(\omega t - \varphi)$ com $R=\sqrt{A^2+B^2}>0$ e fase $\varphi$
adequada. A condi\c{c}\~ao $u(t+2\pi)=u(t)$ para todo $t$ significa
\[
  R\cos\big(\omega t + 2\pi\omega - \varphi\big) = R\cos(\omega t - \varphi)
  \qquad\text{para todo }t .
\]
Como $R\neq 0$, a fun\c{c}\~ao $s\mapsto \cos(s-\varphi)$ \'e $2\pi$-peri\'odica
e n\~ao constante; duas tais fun\c{c}\~oes coincidem para todo $t$ se, e
somente se, seus argumentos diferem por um m\'ultiplo inteiro de $2\pi$, isto
\'e, $2\pi\omega = 2\pi m$ para algum $m\in\Z$. Logo $\omega = m\in\Z$, e como
$\omega>0$, temos $\omega=k$ com $k\in\{1,2,3,\dots\}$.
\end{proof}

\begin{pergunta}\label{perg:inteira}
\textbf{Por que a frequ\^encia \emph{tem} de ser inteira?}
Porque o dom\'\i nio \'e um c\'\i rculo (a reta com $t$ e $t+2\pi$
identificados), e n\~ao um intervalo. Uma onda $\cos(\omega t)$ s\'o ``fecha''
suavemente ao dar uma volta no c\'\i rculo se couber um n\'umero \emph{inteiro}
de oscila\c{c}\~oes em $[0,2\pi]$. Frequ\^encias n\~ao inteiras produziriam uma
descontinuidade ao reencontrar o ponto de partida, violando a periodicidade.
Esta quantiza\c{c}\~ao da frequ\^encia \'e o eco, na linguagem de EDOs, do fato
de que os autovalores do operador $-\d^2/\d t^2$ no c\'\i rculo s\~ao os
quadrados perfeitos $k^2$.
\end{pergunta}

Pelo Lema~\ref{lema:freq}, todo candidato a extremal tem a forma
\begin{equation}\label{eq:harmonico}
  u(t) = A\cos(kt) + B\sin(kt),
  \qquad k\in\{1,2,3,\dots\},\ (A,B)\neq(0,0) .
\end{equation}
Note que $k=0$ daria $u$ constante, exclu\'\i da pela m\'edia nula; e que
qualquer fun\c{c}\~ao da forma \eqref{eq:harmonico} j\'a tem m\'edia nula
automaticamente, pois $\int_0^{2\pi}\cos(kt)\,\d t = \int_0^{2\pi}\sin(kt)\,\d
t = 0$ para $k\ge 1$. Assim \eqref{eq:harmonico} de fato pertence a $\Adm$.

\subsection{O quociente de Rayleigh dos harm\^onicos}

Calculemos $\Q$ na fam\'\i lia \eqref{eq:harmonico}. Usando as integrais
elementares
\[
  \int_0^{2\pi}\cos^2(kt)\,\d t = \int_0^{2\pi}\sin^2(kt)\,\d t = \pi,
  \qquad
  \int_0^{2\pi}\cos(kt)\sin(kt)\,\d t = 0
  \quad (k\ge 1),
\]
obtemos, para $u(t)=A\cos kt + B\sin kt$,
\[
  \int_0^{2\pi} u^2 \,\d t = \pi\,(A^2+B^2) .
\]
Como $u'(t) = -Ak\sin(kt) + Bk\cos(kt)$, o mesmo c\'alculo d\'a
\[
  \int_0^{2\pi} (u')^2 \,\d t = k^2\,\pi\,(A^2+B^2) .
\]
Portanto
\begin{equation}\label{eq:Q-harmonico}
  \Q[u] = \frac{\pi(A^2+B^2)}{k^2\,\pi(A^2+B^2)} = \frac{1}{k^2} .
\end{equation}

%=====================================================================
\section{Conclus\~ao: a constante \'otima e o caso de igualdade}
\label{sec:conclusao}
%=====================================================================

Reunimos agora as pe\c{c}as.

\begin{proof}[Demonstra\c{c}\~ao do Teorema~\ref{thm:wirtinger}]
Pela hip\'otese de exist\^encia (Observa\c{c}\~ao~\ref{obs:existencia}), seja
$u$ um maximizador de $\Q$, normalizado por \eqref{eq:normaliza}. As
Se\c{c}\~oes~\ref{sec:EL}--\ref{sec:edo} mostraram que $u$ resolve
\eqref{eq:oscilador} e, por periodicidade (Lema~\ref{lema:freq}), tem a forma
\eqref{eq:harmonico} para algum inteiro $k\ge 1$. Por \eqref{eq:Q-harmonico},
\[
  \sup_{w\in\Adm\setminus\{0\}}\Q[w] \;=\; \Q[u] \;=\; \frac{1}{k^2} .
\]
Mas o pr\'oprio primeiro harm\^onico $w_1(t)=\cos t$ pertence a
$\Adm\setminus\{0\}$ e tem $\Q[w_1]=1$ por \eqref{eq:Q-harmonico} com $k=1$.
Como $u$ realiza o supremo, devemos ter $1/k^2 = \Q[u] \ge \Q[w_1] = 1$, o que
s\'o \'e poss\'\i vel se $k=1$. Logo
\[
  \sup_{w\in\Adm\setminus\{0\}}\Q[w] = 1 .
\]
Em particular, $\Q[w]\le 1$ para toda $w\in\Adm\setminus\{0\}$, isto \'e,
\[
  \int_0^{2\pi} w^2 \;\le\; \int_0^{2\pi} (w')^2
  \qquad\text{para toda }w\in\Adm,
\]
o que tamb\'em vale trivialmente para $w\equiv 0$. Isto prova a desigualdade.

\medskip
\noindent\emph{Caso de igualdade.}\quad
Se $w\in\Adm$ realiza a igualdade $\int w^2=\int(w')^2$ com $w\not\equiv 0$,
ent\~ao $\Q[w]=1$, ou seja, $w$ \'e um maximizador. Repetindo todo o
argumento das Se\c{c}\~oes~\ref{sec:EL}--\ref{sec:edo} com $w$ no lugar de $u$
(a normaliza\c{c}\~ao foi apenas uma convenei\^encia de escala e n\~ao afeta a
forma da solu\c{c}\~ao), conclu\'\i mos que $w(t)=A\cos kt + B\sin kt$ com
$\Q[w]=1/k^2=1$, donde $k=1$ e
\[
  w(t) = a\cos t + b\sin t .
\]
Reciprocamente, toda fun\c{c}\~ao $w(t)=a\cos t+b\sin t$ pertence a $\Adm$ e
satisfaz $\Q[w]=1$ por \eqref{eq:Q-harmonico}, logo realiza a igualdade. Isto
caracteriza completamente o caso de igualdade.
\end{proof}

\begin{pergunta}\label{perg:constante}
\textbf{Qual passo produz a constante \'otima?}
O passo decisivo \'e \eqref{eq:Q-harmonico} seguido da escolha $k=1$. A
an\'alise da EDO restringe os candidatos a extremal \`a fam\'\i lia discreta
$\{A\cos kt + B\sin kt\}_{k\ge1}$, e nessa fam\'\i lia o quociente vale
exatamente $1/k^2$. O \emph{maior} valor poss\'\i vel ocorre na menor
frequ\^encia admiss\'\i vel, $k=1$, dando $\Q=1$. \'E literalmente a
escolha do primeiro harm\^onico que fixa a constante \'otima da
desigualdade.
\end{pergunta}

\begin{observacao}[Leitura espectral]
A conta acima \'e a vers\~ao ``por EDO'' do fato de que os autovalores do
problema $-u''=\nu u$ com condi\c{c}\~oes peri\'odicas em $[0,2\pi]$ s\~ao
$\nu_k=k^2$ ($k=0,1,2,\dots$), com autofun\c{c}\~oes $1,\cos t,\sin t,\cos
2t,\sin 2t,\dots$. A m\'edia nula remove o autovalor $\nu_0=0$ (a
constante), e o menor autovalor restante, $\nu_1=1$, \'e exatamente a constante
\'otima de Wirtinger.
\end{observacao}

%=====================================================================
\section{A forma reescalada (per\'\i odo $L$)}
\label{sec:escala}
%=====================================================================

A escolha do per\'\i odo $2\pi$ foi apenas uma normaliza\c{c}\~ao. Eis a forma
geral, obtida por uma simples mudan\c{c}a de vari\'avel.

\begin{corolario}[Wirtinger com per\'\i odo $L$]\label{cor:escala}
Seja $u\colon\R\to\R$ de classe $C^1$, $L$-peri\'odica e de m\'edia nula,
$\int_0^L u(x)\,\d x = 0$. Ent\~ao
\[
  \int_0^L u(x)^2 \,\d x \;\le\; \frac{L^2}{4\pi^2}\int_0^L u'(x)^2 \,\d x ,
\]
com igualdade se, e somente se,
$u(x)=a\cos\!\big(\tfrac{2\pi x}{L}\big)+b\sin\!\big(\tfrac{2\pi x}{L}\big)$.
\end{corolario}

\begin{proof}
Reparametrizamos para o per\'\i odo $2\pi$. Defina $w\colon\R\to\R$ por
\[
  w(t) = u\!\left(\frac{L}{2\pi}\,t\right) .
\]
Ent\~ao $w$ \'e $C^1$ e $2\pi$-peri\'odica, pois
$w(t+2\pi)=u\big(\tfrac{L}{2\pi}t + L\big)=u\big(\tfrac{L}{2\pi}t\big)=w(t)$.
Com a substitui\c{c}\~ao $x=\tfrac{L}{2\pi}t$, $\d x=\tfrac{L}{2\pi}\,\d t$,
a m\'edia de $w$ \'e
\[
  \int_0^{2\pi} w(t)\,\d t
  = \frac{2\pi}{L}\int_0^L u(x)\,\d x = 0 ,
\]
logo $w\in\Adm$. Pela regra da cadeia, $w'(t)=\tfrac{L}{2\pi}\,u'\!\big(\tfrac{L}{2\pi}t\big)$.
Convertendo as duas integrais de Wirtinger para a vari\'avel $x$:
\[
  \int_0^{2\pi} w(t)^2\,\d t
  = \frac{2\pi}{L}\int_0^L u(x)^2\,\d x,
\]
\[
  \int_0^{2\pi} w'(t)^2\,\d t
  = \left(\frac{L}{2\pi}\right)^{\!2}\frac{2\pi}{L}\int_0^L u'(x)^2\,\d x
  = \frac{L}{2\pi}\int_0^L u'(x)^2\,\d x .
\]
Aplicando o Teorema~\ref{thm:wirtinger} a $w$, isto \'e
$\int_0^{2\pi}w^2\le\int_0^{2\pi}(w')^2$, obtemos
\[
  \frac{2\pi}{L}\int_0^L u^2 \;\le\; \frac{L}{2\pi}\int_0^L (u')^2 ,
\]
e multiplicando por $\tfrac{L}{2\pi}$ chega-se a
$\int_0^L u^2 \le \tfrac{L^2}{4\pi^2}\int_0^L (u')^2$. O caso de igualdade
corresponde, via a mesma substitui\c{c}\~ao, a $w(t)=a\cos t+b\sin t$, ou seja
$u(x)=a\cos\!\big(\tfrac{2\pi x}{L}\big)+b\sin\!\big(\tfrac{2\pi x}{L}\big)$.
\end{proof}

%=====================================================================
\section{Aplica\c{c}\~ao: a desigualdade isoperim\'etrica no plano}
\label{sec:isop}
%=====================================================================

A desigualdade isoperim\'etrica afirma que, dentre todas as curvas fechadas
simples de um dado comprimento, o c\'\i rculo encerra a maior \'area. Daremos
a elegante prova de Hurwitz, que reduz tudo \`a desigualdade de Wirtinger.

\begin{teorema}[Desigualdade isoperim\'etrica]\label{thm:isop}
Seja $\gamma$ uma curva plana fechada simples, de classe $C^1$, com
comprimento $L$ e \'area encerrada $A$. Ent\~ao
\[
  4\pi A \;\le\; L^2 ,
\]
com igualdade se, e somente se, $\gamma$ \'e um c\'\i rculo.
\end{teorema}

\subsection{Prepara\c{c}\~ao: parametriza\c{c}\~ao por velocidade constante}

Parametrizamos $\gamma$ no intervalo $[0,2\pi]$ proporcionalmente ao
comprimento de arco, escrevendo $\gamma(t)=(f(t),g(t))$ com $f,g$ de classe
$C^1$ e $2\pi$-peri\'odicas. ``Velocidade constante'' significa que
$|\gamma'(t)|$ \'e constante. Como o comprimento total \'e
$L=\int_0^{2\pi}|\gamma'(t)|\,\d t$, essa constante vale $L/2\pi$, isto \'e,
\begin{equation}\label{eq:velocidade}
  f'(t)^2 + g'(t)^2 = \left(\frac{L}{2\pi}\right)^{\!2}
  = \frac{L^2}{4\pi^2}
  \qquad\text{para todo }t .
\end{equation}
Integrando \eqref{eq:velocidade} sobre $[0,2\pi]$ obtemos a identidade que
usaremos ao final:
\begin{equation}\label{eq:int-velocidade}
  \int_0^{2\pi}\!\big(f'^2 + g'^2\big)\,\d t
  = 2\pi\cdot\frac{L^2}{4\pi^2}
  = \frac{L^2}{2\pi} .
\end{equation}
Finalmente, \emph{transladamos} a curva na dire\c{c}\~ao horizontal de modo que
\begin{equation}\label{eq:f-media-zero}
  \int_0^{2\pi} f(t)\,\d t = 0 .
\end{equation}
Isto \'e leg\'\i timo: substituir $f$ por $f-c$ (com $c$ a m\'edia de $f$) \'e
uma transla\c{c}\~ao r\'\i gida da curva, que n\~ao altera nem a \'area $A$ nem
o comprimento $L$; e n\~ao altera $f'$, de modo que \eqref{eq:velocidade}
continua valendo. Ap\'os a transla\c{c}\~ao, $f\in\Adm$.

\subsection{A \'area via Teorema de Green}

\begin{lema}[F\'ormula da \'area]\label{lema:area}
Com a parametriza\c{c}\~ao acima (orientada positivamente),
\[
  A = \int_0^{2\pi} f(t)\,g'(t)\,\d t .
\]
\end{lema}

\begin{proof}
Pelo Teorema de Green, a \'area encerrada por uma curva fechada simples
positivamente orientada \'e
\[
  A = \frac12\oint_\gamma (x\,\d y - y\,\d x)
    = \frac12\int_0^{2\pi}\big(f g' - g f'\big)\,\d t .
\]
Integrando $\int_0^{2\pi} f g'\,\d t$ por partes, com termo de fronteira nulo
pela periodicidade (como em \eqref{eq:partes}),
\[
  \int_0^{2\pi} f\,g'\,\d t
  = \big[\,f g\,\big]_0^{2\pi} - \int_0^{2\pi} f'\,g\,\d t
  = - \int_0^{2\pi} f'\,g\,\d t .
\]
Logo $-\int g f' = \int f g'$, e substituindo,
\[
  A = \frac12\Big(\int f g' - \int g f'\Big)
    = \frac12\Big(\int f g' + \int f g'\Big)
    = \int_0^{2\pi} f\,g'\,\d t . \qedhere
\]
\end{proof}

\subsection{O n\'ucleo do argumento}

Aqui se combina tudo. Partimos de $2A=2\int fg'$ e usamos o ``truque do
quadrado'': para quaisquer n\'umeros reais $f$ e $g'$,
\begin{equation}\label{eq:truque}
  2fg' = f^2 + g'^2 - (f - g')^2 .
\end{equation}
Integrando \eqref{eq:truque} e descartando o termo $-\int (f-g')^2\le 0$,
\begin{equation}\label{eq:2A}
  2A = \int_0^{2\pi} 2 f g'
     = \int_0^{2\pi}\!\big(f^2 + g'^2\big) - \int_0^{2\pi}(f-g')^2
     \;\le\; \int_0^{2\pi} f^2 + \int_0^{2\pi} g'^2 .
\end{equation}
Agora aplicamos a desigualdade de Wirtinger (Teorema~\ref{thm:wirtinger}) \`a
fun\c{c}\~ao $f\in\Adm$, que tem m\'edia nula por \eqref{eq:f-media-zero}:
\begin{equation}\label{eq:wirtinger-f}
  \int_0^{2\pi} f^2 \;\le\; \int_0^{2\pi} (f')^2 .
\end{equation}
Substituindo \eqref{eq:wirtinger-f} em \eqref{eq:2A},
\[
  2A \;\le\; \int_0^{2\pi}(f')^2 + \int_0^{2\pi}(g')^2
      = \int_0^{2\pi}\!\big(f'^2 + g'^2\big)
      \overset{\eqref{eq:int-velocidade}}{=} \frac{L^2}{2\pi} .
\]
Portanto $2A\le L^2/2\pi$, isto \'e,
\[
  \boxed{\,4\pi A \le L^2\,}.
\]
Isto prova a desigualdade isoperim\'etrica.

\subsection{O caso de igualdade: a curva \'e um c\'\i rculo}

Suponha agora que vale a igualdade $4\pi A = L^2$. Ent\~ao as duas
desigualdades que usamos --- \eqref{eq:2A} e \eqref{eq:wirtinger-f} --- devem
ser, ambas, igualdades.

\emph{Igualdade em \eqref{eq:2A}} significa $\int_0^{2\pi}(f-g')^2 = 0$.
Como o integrando \'e cont\'\i nuo e n\~ao negativo, conclu\'\i mos
\begin{equation}\label{eq:g'=f}
  g'(t) = f(t) \qquad\text{para todo }t .
\end{equation}

\emph{Igualdade em \eqref{eq:wirtinger-f}} \'e a igualdade na desigualdade de
Wirtinger; pelo Teorema~\ref{thm:wirtinger}, isso for\c{c}a
\begin{equation}\label{eq:f-harmonico}
  f(t) = a\cos t + b\sin t
\end{equation}
para certas constantes $a,b$. Por \eqref{eq:g'=f}, $g'=f=a\cos t+b\sin t$,
e integrando,
\begin{equation}\label{eq:g}
  g(t) = a\sin t - b\cos t + c
\end{equation}
para alguma constante $c$. Examinemos a curva resultante. Transladando-a
verticalmente por $-c$ (o que n\~ao muda sua forma), obtemos
$\gamma(t)-(0,c) = \big(a\cos t + b\sin t,\; a\sin t - b\cos t\big)$, e
\[
  f(t)^2 + \big(g(t)-c\big)^2
  = (a\cos t + b\sin t)^2 + (a\sin t - b\cos t)^2
  = a^2 + b^2 ,
\]
constante. Logo $\gamma$ \'e um c\'\i rculo de centro $(0,c)$ e raio
$r=\sqrt{a^2+b^2}$.

Reciprocamente, para um c\'\i rculo de raio $r$ tem-se $L=2\pi r$ e
$A=\pi r^2$, de modo que $4\pi A = 4\pi\cdot\pi r^2 = (2\pi r)^2 = L^2$: o
c\'\i rculo realiza a igualdade. Isto encerra a prova do
Teorema~\ref{thm:isop}.
\hfill$\qed$

\begin{observacao}[Consist\^encia com a velocidade constante]
Vale a pena verificar que a curva extremal \eqref{eq:f-harmonico}--\eqref{eq:g}
de fato tem velocidade constante, como exigido em \eqref{eq:velocidade}. Com
$f'=-a\sin t+b\cos t$ e $g'=f=a\cos t+b\sin t$,
\[
  f'^2 + g'^2
  = (-a\sin t+b\cos t)^2 + (a\cos t+b\sin t)^2 = a^2+b^2 = r^2 ,
\]
constante, e igual a $L^2/4\pi^2$ pois $L=2\pi r$. Tudo \'e consistente: o
c\'\i rculo percorrido a velocidade angular constante \'e exatamente a curva
que satura simultaneamente as duas desigualdades.
\end{observacao}

%=====================================================================
\section{S\'\i ntese e exerc\'\i cios}
%=====================================================================

\subsection*{Roteiro l\'ogico}
O argumento inteiro pode ser lido como uma cadeia de implica\c{c}\~oes
t\'\i picas de EDO:
\[
  \begin{gathered}
  \text{maximizador de }\Q
  \;\longrightarrow\;
  \text{Euler--Lagrange } u''+\omega^2u=0
  \;\longrightarrow\;
  u = A\cos kt + B\sin kt \\[2pt]
  \longrightarrow\quad
  \Q=\tfrac{1}{k^2}\le 1 .
  \end{gathered}
\]
A m\'edia nula remove a frequ\^encia $k=0$
(Pergunta~\ref{perg:media}); a periodicidade faz o termo de fronteira
desaparecer (Pergunta~\ref{perg:fronteira}) e quantiza a frequ\^encia
(Pergunta~\ref{perg:inteira}); e a escolha $k=1$ produz a constante \'otima
(Pergunta~\ref{perg:constante}). A desigualdade isoperim\'etrica segue por um
\'unico ``truque do quadrado'' aliado a Wirtinger.

\subsection*{Exerc\'\i cios}

\begin{exercicio}
Mostre que, sem a hip\'otese de m\'edia nula, vale a desigualdade de
Poincar\'e $\int_0^{2\pi}\big(u-\bar u\big)^2\le \int_0^{2\pi}(u')^2$, onde
$\bar u=\frac{1}{2\pi}\int_0^{2\pi}u$ \'e a m\'edia de $u$. (Sugest\~ao:
aplique o Teorema~\ref{thm:wirtinger} a $u-\bar u$.)
\end{exercicio}

\begin{exercicio}
Verifique diretamente, sem variar nada, que $\int_0^{2\pi}u^2\le
\int_0^{2\pi}(u')^2$ para $u(t)=\sin t + \tfrac13\sin 3t$. Compare os dois
lados e identifique a contribui\c{c}\~ao de cada harm\^onico.
\end{exercicio}

\begin{exercicio}
Refa\c{c}a a dedu\c{c}\~ao da equa\c{c}\~ao de Euler--Lagrange para o problema
de Dirichlet: maximizar $\int_0^{\pi}u^2$ sujeito a $\int_0^{\pi}(u')^2=1$
\emph{e} $u(0)=u(\pi)=0$ (em vez de periodicidade). Mostre que agora as
fun\c{c}\~oes de teste devem satisfazer $v(0)=v(\pi)=0$, e que os extremais
s\~ao $u(t)=\sin(kt)$. Conclua a desigualdade $\int_0^\pi u^2\le\int_0^\pi
(u')^2$ com igualdade em $u=\sin t$.
\end{exercicio}

\begin{exercicio}
Prove o Lema~\ref{lema:freq} de um modo alternativo, impondo diretamente
$u(0)=u(2\pi)$ e $u'(0)=u'(2\pi)$ na solu\c{c}\~ao geral
$A\cos(\omega t)+B\sin(\omega t)$ e mostrando que um sistema com solu\c{c}\~ao
$(A,B)\neq(0,0)$ exige $\cos(2\pi\omega)=1$.
\end{exercicio}

\begin{exercicio}
Na prova isoperim\'etrica, explique por que foi essencial parametrizar a curva
com velocidade constante. O que falha no passo \eqref{eq:int-velocidade} se a
parametriza\c{c}\~ao for arbitr\'aria? (Sugest\~ao: a desigualdade
$4\pi A\le L^2$ continua verdadeira, mas a identidade
$\int(f'^2+g'^2)=L^2/2\pi$ usa crucialmente $|\gamma'|\equiv L/2\pi$.)
\end{exercicio}

\subsection*{Refer\^encias e coment\'arios}
A prova variacional aqui apresentada \'e uma reformula\c{c}\~ao, em linguagem
de EDO, da prova cl\'assica via s\'eries de Fourier (que aparece, por exemplo,
em textos de an\'alise de Fourier). A aplica\c{c}\~ao \`a desigualdade
isoperim\'etrica \'e devida a A.~Hurwitz (1901). Para a justificativa
rigorosa da exist\^encia do maximizador (Observa\c{c}\~ao~\ref{obs:existencia})
e da regularidade, consulte um tratamento do m\'etodo direto do c\'alculo de
varia\c{c}\~oes em espa\c{c}os de Sobolev.

\end{document}
